% =====================================================================
%  Conferenze, GAMM Roma (Checkpoint Festival, 3a edizione) — Panel 1
%  Non solo intrattenimento: videogiochi, società e tutela dei minori
%  Fonte: trascrizione automatica "Track 1" di un panel tenuto al GAMM
%  (Game Museum) di Roma durante la terza edizione del Checkpoint
%  Festival. Partecipanti: una moderatrice (non nominata), la presidente
%  del Corecom Lazio, un game designer di Milestone e un rappresentante
%  dell'associazione di categoria dell'industria italiana. La
%  registrazione comincia a panel già avviato e si interrompe durante i
%  saluti finali.
%  Etichette con prefisso "c1-1:". Nessun pacchetto aggiuntivo.
% =====================================================================

\chapter[Videogiochi, società e tutela dei minori]{Non solo intrattenimento:
videogiochi, società e tutela dei minori}
\label{cap:c1-1-tutela-minori}

Il primo intervento registrato al GAMM è un \textbf{panel a più voci} sul
rapporto tra videogiochi e società, con un taglio preciso: la
\textbf{tutela dei minori} negli ambienti di gioco, online e offline. La
registrazione comincia a frase già iniziata, mentre la moderatrice sta
spiegando la tesi del panel: i videogiochi ``sono intrattenimento, è giusto che
lo siano, ma non sono solo quello. Sono un'industria, sono un settore in
continuo sviluppo e che sta piano piano arrivando a influenzare in modo
positivo anche la società, la cultura e la comunicazione digitale''.

Al tavolo siedono tre ospiti, presentati dalla moderatrice\footnote{La
moderatrice non dice il proprio nome nella parte registrata. Dal contesto
(consegna dei premi agli ospiti ``in questi giorni'', ringraziamenti ``a
casa'' e a chi segue online) fa parte dell'organizzazione del festival.}:
\begin{itemize}
  \item \textbf{Eleonora Zazza}\footnote{Il cognome compare come ``Zazza'' nella
        presentazione e come ``Zarza'' durante la premiazione finale; si usa la
        prima forma.}, avvocata e \textbf{presidente del Corecom Lazio}, il
        Comitato Regionale per le Comunicazioni\footnote{La trascrizione rende la
        sigla in vari modi (``Corecom'', ``Coreo Palazzo'', ``Cori Con Lazio''):
        si tratta sempre del Co.Re.Com.\ del Lazio, organo regionale che opera
        anche come organo funzionale dell'Autorità per le garanzie nelle
        comunicazioni.}, che secondo la moderatrice sta portando avanti ``due
        bellissime iniziative'' per il programma 2027;
  \item \textbf{Federico Spada}\footnote{Il ruolo è trascritto come ``candidato
        di Milestone'', con ogni probabilità una storpiatura (game designer o
        \emph{creative director}); più avanti la moderatrice lo chiama anche
        ``Dirigo'', certamente un errore per ``Federico''. Il cognome non è
        verificabile e si riporta come trascritto.}, di \textbf{Milestone}, la
        casa di produzione italiana nota per i giochi di corse (\emph{MotoGP},
        \emph{RIDE}\footnote{``La PvS'' nella trascrizione: con ogni
        probabilità una delle altre serie motoristiche di Milestone; la lettura
        \emph{RIDE} è un'ipotesi.}), che ha contribuito allo sviluppo
        dell'ultimo titolo dello studio, uscito a settembre\footnote{La
        trascrizione dice ``il Christopher Flash'' e poi ``Scrivere'': quasi
        certamente \emph{Screamer}, il gioco di corse arcade di Milestone,
        ripreso come titolo sperimentale rispetto alle serie su licenza dello
        studio.};
  \item \textbf{Davide}\footnote{Nella presentazione il nome appare come
        ``Davide Mancini'', ma il passaggio è molto disturbato; l'appartenenza
        all'associazione di categoria dell'industria italiana si ricava da ciò
        che dice lui stesso (``noi come associazione siamo parte del PEGI, [\dots]
        di due federazioni europee''). È definito ``direttore artistico'' di un
        evento (trascritto ``The Women's World'') e scherza sul fatto che sia
        ``uno step per arrivare a Sanremo'': si tratta verosimilmente della
        cerimonia di premiazione dei videogiochi italiani organizzata
        dall'associazione. Sia il cognome sia il nome dell'evento restano
        incerti.}, che fa ``da ponte'' tra tutti i temi del panel: ha un passato
        nel \textbf{giornalismo videoludico}, lavora ``da quasi quindici anni''
        nell'industria e oggi opera nell'associazione di categoria.
\end{itemize}

Il capitolo segue l'ordine del panel:
\begin{enumerate}
  \item l'intervento d'apertura della presidente del Corecom
        (sezione~\ref{sec:c1-1-zazza});
  \item il suo racconto da genitore e la questione dei ``nativi digitali''
        (sezione~\ref{sec:c1-1-genitore});
  \item il ``gioco'' proposto dalla moderatrice: da dove partireste per
        tutelare i minori? Le risposte dell'associazione
        (sezione~\ref{sec:c1-1-autoregolazione}) e dello sviluppatore
        (sezione~\ref{sec:c1-1-sviluppatore});
  \item l'evoluzione dei videogiochi italiani
        (sezione~\ref{sec:c1-1-italia});
  \item che cosa rende riconoscibile un gioco Milestone
        (sezione~\ref{sec:c1-1-dettaglio});
  \item il nodo del \textbf{contatto online} e come lo si affronta in fase di
        progettazione (sezione~\ref{sec:c1-1-online});
  \item la chiusura e la premiazione (sezione~\ref{sec:c1-1-chiusura}).
\end{enumerate}

% ---------------------------------------------------------------------
\section{L'apertura del Corecom: una responsabilità condivisa}
\label{sec:c1-1-zazza}

Eleonora Zazza prende la parola definendo il momento ``storico'': è al GAMM
come presidente del Corecom Lazio ma anche, aggiunge, come
\textbf{madre}\footnote{La frase sui figli è quasi incomprensibile nella
trascrizione (``a me non mi dà due ragazze di cui uno a Francia\dots''); più
avanti dice chiaramente di avere un figlio adolescente di quattordici anni che
``approccia'' i videogiochi.}, con un figlio di quattordici anni. Si discute
molto, dice, della tutela dei minori, ``ed è giusto'': la tutela dei minori
nell'uso della comunicazione e nel mondo dei videogiochi è la missione del
comitato ``sia a livello istituzionale, sia a livello di madre''.

Ringrazia il museo (il GAMM), i rappresentanti dell'industria presenti e il
pubblico, perché la loro presenza è ``un segnale di attenzione, di
collaborazione'': la tutela dei minori nei videogiochi è una responsabilità che
\textbf{deve diventare un impegno condiviso} da tutte le parti in causa:
\begin{itemize}
  \item le \textbf{autorità};
  \item l'\textbf{industria};
  \item i \textbf{genitori};
  \item e, come ultimo riferimento, la \textbf{scuola}.
\end{itemize}
Sottolinea che l'attenzione richiesta è ``diversa'': non all'industria in sé,
ma agli \textbf{effetti dell'industria sulla società}.

\subsection{Opportunità e rischi}

La tutela dei minori in ambiente digitale è per lei una responsabilità
istituzionale che richiede tre cose: \textbf{attenzione costante},
\textbf{capacità di prevenzione} e \textbf{collaborazione}, ``soprattutto''
in ottica collaborativa. I videogiochi sono oggi una realtà importante nella
vita di bambini e adolescenti e offrono opportunità di:
\begin{itemize}
  \item \textbf{apprendimento} (vengono usati anche in situazioni di
        difficoltà di apprendimento);
  \item \textbf{creatività};
  \item \textbf{socializzazione}.
\end{itemize}
Ma pongono anche questioni ``che non possiamo sottovalutare''. La presidente
precisa di non voler mettere in discussione il valore del settore né alimentare
una \textbf{contrapposizione tra tecnologia e sicurezza}. Al contrario: proprio
un'industria capace di raggiungere livelli così avanzati di innovazione deve
essere \textbf{protagonista} nella costruzione di strumenti sempre più efficaci
per proteggere gli utenti più giovani, perché ``la capacità di innovare deve
accompagnarsi alla capacità di assumersi delle responsabilità''. E cita: ``un
grande potere deve essere sempre collegato a grandi responsabilità''.

\subsection{I dati}

Due fonti, per misurare la portata della sfida:
\begin{description}
  \item[Associazione Meter.] Una ricerca condotta su \textbf{467 bambini tra i
        9 e i 15 anni} che usano \emph{Roblox}\footnote{Trascritto ``il
        ROCROC''; il riferimento a \emph{Roblox}, piattaforma molto diffusa tra
        i più giovani, è un'interpretazione ragionevole ma non certa.}:
        \begin{itemize}
          \item il \textbf{70\%} ha dichiarato di aver incontrato almeno una
                situazione di rischio durante il gioco;
          \item il \textbf{45\%} ha riferito di essere stato esposto al
                pericolo dell'\textbf{adescamento} da parte di
                sconosciuti\footnote{La trascrizione ripete due volte la stessa
                frase (``almeno una situazione di rischio'') con le due
                percentuali e poi dice ``al pericolo del traduriggio,
                adescamento'': si tratta quasi certamente del \emph{grooming},
                cioè l'adescamento online. L'abbinamento del 45\% a questo
                rischio è una ricostruzione.};
          \item solo il \textbf{10\%} ha mostrato piena consapevolezza della
                gravità del fenomeno.
        \end{itemize}
        Sono dati riferiti a uno specifico campione, ``che non devono essere
        generalizzati'', ma che sollevano interrogativi ``ai quali non possiamo
        sottrarci''.
  \item[Telefono Azzurro.] Nel dossier presentato nel 2026 si evidenzia che,
        nel corso del 2025, il \textbf{7,2\%} dei casi gestiti dal servizio
        \textbf{114 Emergenza Infanzia} riguardava situazioni riconducibili ad
        abuso e sfruttamento sessuale, adescamento, ricatto sessuale e
        materiale di abuso sessuale su minori. Il dato riguarda fenomeni più
        ampi del videogioco, ma richiama la necessità di rendere più sicuri
        \textbf{tutti gli ambienti digitali frequentati} dai ragazzi.
\end{description}

Il principio che la presidente ripete ``sempre'' è che \textbf{mondo offline e
mondo online non sono separati}: come si parla di sicurezza nelle piazze, si
deve parlare di sicurezza negli ambienti digitali. ``Non stiamo parlando di
fantascienza, stiamo parlando di realtà.''

\subsection{L'appello all'industria}

Non basta più, dice, chiedere alle famiglie, o \emph{solo} alle famiglie, di
vigilare: bisogna \textbf{metterle nelle condizioni} di proteggere davvero i
ragazzi, e questo richiede un'assunzione di responsabilità da parte
dell'industria. Rivolgendosi ai rappresentanti delle imprese, li invita a fare
della tutela dei minori \textbf{un elemento centrale delle scelte progettuali e
produttive}, e in particolare a investire in:
\begin{enumerate}
  \item \textbf{sistemi di protezione adeguati all'età};
  \item \textbf{strumenti di controllo parentale} semplici ed efficaci;
  \item \textbf{modalità di comunicazione online} più sicure;
  \item maggiore \textbf{trasparenza nei meccanismi commerciali};
  \item attenzione anche ai \textbf{videogiochi offline}, perché la tutela non
        riguarda soltanto i pericoli della rete, ma anche l'\textbf{adeguatezza
        dei contenuti} e l'\textbf{equilibrio nell'utilizzo}.
\end{enumerate}

Riconosce che il settore ha già fatto ``passi importanti'': i \textbf{sistemi di
classificazione per età} e gli strumenti di protezione esistenti. Ma pone tre
domande: sono \textbf{sufficienti}? Sono \textbf{realmente conosciuti}? (``È qui
che dobbiamo aiutare, perché magari ci sono ma non sono conosciuti.'')
Rispondono davvero \textbf{alle esigenze delle famiglie}? La sicurezza non può
essere soltanto un'opzione offerta all'utente\footnote{La trascrizione dice
``una possibilità offensiva'': forse ``una possibilità offerta'' o
``un'opzione''.}: deve diventare un \textbf{criterio primario nella progettazione
dell'esperienza di gioco}.

Il Corecom, nell'ambito delle sue competenze, si occupa di \textbf{educazione
digitale}, prevenzione del \textbf{cyberbullismo} e promozione di un uso
responsabile delle tecnologie; ma oggi serve rafforzare il dialogo con chi
quelle tecnologie le sviluppa e le porta sul mercato. Di qui l'invito preciso:
\begin{quote}
  promuoviamo un confronto continuo, costante, strutturato tra industrie,
  istituzioni e famiglie per individuare impegni concreti, verificabili e
  orientabili per la protezione dei minori.
\end{quote}
Un confronto che non si esaurisca nelle dichiarazioni di principio ma produca
iniziative condivise: \textbf{formazione per le famiglie}, \textbf{corsi di
educazione digitale}, \textbf{soluzioni tecniche} capaci di prevenire i rischi.
Non si tratta di imporre una visione contraria allo sviluppo del settore, ma di
riconoscere che \textbf{la qualità di un prodotto tecnologico si misura anche
dalla sua capacità di rispettare i diritti di chi lo utilizza}.

La considerazione finale che ``affida al tavolo'': chi sviluppa videogiochi ha
le competenze per immaginare e realizzare mondi virtuali sempre più sofisticati;
può anche fare in modo che siano \textbf{luoghi sempre più sicuri}. ``Non
possiamo pretendere che un bambino abbia gli strumenti per riconoscere ogni
pericolo, e neanche che i genitori ce l'abbiano''; possiamo invece pretendere un
impegno serio da parte degli adulti e di chi progetta gli ambienti. La tutela
dei minori, conclude, \textbf{non è un limite all'innovazione ma una misura
della sua qualità}\footnote{``Non deve essere considerata un sacco
all'innovazione'': quasi certamente ``un freno'' o ``un ostacolo''.}: il futuro
dell'industria ``non deve essere soltanto più innovativo, deve essere anche più
sicuro, più consapevole e più attento ai diritti dei bambini''.

L'impegno concreto preso dal Corecom Lazio è portare \textbf{nelle scuole} e
all'attenzione dei \textbf{genitori} il tema dell'``approccio al giocare bene''.
Il videogioco, nella sua immagine, deve poter essere ciò che un tempo era
\textbf{la piazza}: un luogo che i ragazzi abitano e che devono essere messi in
condizione di abitare. È un mondo nuovo e, come tutti i mondi nuovi, va
affrontato con l'\textbf{alfabetizzazione}.

% ---------------------------------------------------------------------
\section{Il punto di vista di un genitore}
\label{sec:c1-1-genitore}

La moderatrice riprende l'accenno iniziale al figlio adolescente e chiede alla
presidente se sia anche lei videogiocatrice, o se stia sperimentando i
videogiochi grazie ai figli. La risposta è sincera: \textbf{no, non è una
videogiocatrice}. Però, già prima di diventare presidente del Corecom, quando i
figli hanno cominciato a parlarle di videogiochi, ha cercato di
``inserirsi'' e di capire insieme a loro. Confessa che, come genitore,
\textbf{non sapeva} certe cose sui videogiochi\footnote{La frase è troncata
(``non sapevo che i videogiochi\dots fossero mai lo contesto''). Dal seguito si
capisce che si riferisce all'esistenza delle chat nei giochi e alle loro
implicazioni: ``ho scoperto\dots le chat?''.}, e che come lei ``tanti altri adulti
non lo sanno''. Da lì, prima ancora del suo ruolo istituzionale, si è posta il
problema della \textbf{conoscenza}.

Il ragionamento che ne trae è istituzionale: non bisogna partire dalla propria
esperienza, ma da quella di chi ``magari non ce l'ha'', di chi non ha le
capacità per orientarsi; l'istituzione ``deve guardare a tutti, non a un caso
specifico''. Racconta che il marito è più esperto (``all'epoca i videogiochi non
erano online, ma esistevano'') e che il sabato o la domenica gioca con i figli:
diventa \textbf{un momento di condivisione}. ``Non siamo qui per demonizzare'',
ribadisce, ma per chiamare l'attenzione sulla necessità di tutela e di
responsabilità condivisa; e lo stesso discorso vale per i social e per
qualunque ambiente digitale.

\subsection{Il divieto non basta}

La moderatrice osserva che gli ambienti digitali sono nuovi per molti e ancora
poco regolamentati; cita un recente emendamento del governo sul
\textbf{divieto dei social} per i minori. La presidente si dice d'accordo con il
divieto, perché serve a ``intervenire sull'effetto patogeno'' di un mondo
digitale lasciato diventare ``troppo pervasivo'', come si fece per le sigarette
e per l'alcol. Ma aggiunge subito: \textbf{il divieto non basta}, perché quello
è il mondo che i ragazzi abitano e abiteranno. Serve
\textbf{conoscenza}, \textbf{consapevolezza} e la costruzione di un
\textbf{pensiero critico}.

Racconta di essere nata in un mondo senza internet e senza telefonini (``i primi
erano delle cose enormi, accessibili solo a pochi''): la sua generazione è stata
alfabetizzata al mondo offline, e allo stesso modo bisogna
\textbf{alfabetizzare i giovani al mondo online}\footnote{La trascrizione dice
``affabbrizzati'' e ``affabbrizzare'': ``alfabetizzati'',
``alfabetizzare''.}. Per gli adulti il digitale è stato una rivoluzione; per i
ragazzi no: ci sono nati.

\subsection{Attenzione al ``nativo digitale''}

L'ultima riflessione, ``da genitore'', riguarda il concetto di \textbf{nativo
digitale}. Descrive una realtà oggettiva, ma può essere un'\textbf{arma a doppio
taglio}: rischia di diventare un \textbf{esonero di responsabilità} per l'adulto
(``tu che lo conosci, aiutami tu''), che delega a un ragazzo che magari non è
stato alfabetizzato nel modo corretto. Essere nativi digitali è una condizione
di vita, non una laurea in scienze della comunicazione o in ingegneria: non
significa che il ragazzo sia davvero \textbf{formato} ad affrontare quei rischi.

% ---------------------------------------------------------------------
\section{Da dove partire: l'autoregolamentazione dell'industria}
\label{sec:c1-1-autoregolazione}

La moderatrice propone ``un gioco'' agli altri due ospiti: se doveste creare
una forma di tutela, \textbf{qual è il primo punto su cui andreste a
intervenire?}

Davide comincia dalla storia. L'industria dei videogiochi, dice, ha avuto ``una
fortuna'': pur dovendo ancora fare molta strada nel comunicare correttamente
alcune cose, è un'industria dell'intrattenimento che a un certo punto della sua
storia ha deciso di \textbf{autoregolarsi}. In parte per \textbf{autocoscienza},
in parte per \textbf{motivi commerciali}: se negli Stati Uniti un gioco contiene
contenuti non adatti ai minori, ``non può andare a Walmart'', e questo è un
problema commerciale. Come ogni industria dell'intrattenimento, ammette, è
sempre orientata al profitto; ma questo ha comunque generato una forma di
autocoscienza.

I sistemi di classificazione sono nati \textbf{dall'interno dell'industria}:
\begin{itemize}
  \item prima l'\textbf{ESRB}\footnote{Trascritto ``ISRB'' e poi ``SBR'':
        l'Entertainment Software Rating Board, il sistema di classificazione
        nordamericano.}, americano;
  \item poi il \textbf{PEGI}\footnote{Trascritto anche ``BEGI'': il Pan European
        Game Information, il sistema europeo.}, europeo.
\end{itemize}
Per lui è anche un approccio ``corretto'', perché chiedere a qualcuno di esterno
di regolamentare qualcosa di molto complesso rischia di creare incomprensioni.
Gli strumenti esistono; forse vanno \textbf{potenziati}, sfruttando proprio
questa ``scia dell'autocoscienza''.

\subsection{Gli strumenti per le famiglie}

Segnala, ``non perché è una cosa di categoria'', il sito
\textbf{tuttosuivideogiochi.it}, sistemato ``nelle ultime settimane'', che
racconta l'\textbf{approccio consapevole ai videogiochi}. L'associazione fa
parte del PEGI e di due federazioni europee, e il sito va nella stessa direzione
indicata dalla presidente: comunicare \textbf{a tutti}, non solo a chi gioca,
come gestire il videogioco in ambito familiare. Spiega, per esempio, come
funzionano i \textbf{controlli parentali} sui diversi dispositivi, ciascuno con
le sue regole (``io c'ho la Switch\dots non è banale, non è assolutamente
banale'').

Cita poi un'iniziativa europea ``molto carina'': un \textbf{patto tra genitori e
figli} per stabilire \textbf{regole condivise} sull'uso dei videogiochi. È un
approccio che secondo lui funziona e va ``assolutamente rinforzato''.

\subsection{L'autoregolazione come inizio di un dialogo}

Davide insiste sul termine \emph{autoregolazione}, non per evitare le
regolazioni rigide, ma perché l'autoregolazione è \textbf{l'inizio di un
dialogo}; e comunque si muove dentro un quadro di normative, soprattutto
europee, che regolano l'accesso e l'uso di qualsiasi prodotto digitale. Richiama
la contrapposizione tra \textbf{apocalittici e integrati}\footnote{Il
riferimento, non esplicitato, è al saggio di Umberto Eco \emph{Apocalittici e
integrati} (1964) sulla cultura di massa.}: ``l'unica via che abbiamo è
l'integrazione''. Le domande e i timori sono legittimi; ma l'industria
videoludica, enorme e piena di ``casi d'uso assurdi'', è anche
\textbf{incredibilmente maturata} nel prendere coscienza di sé, pur con ``ancora
tanto margine''. Se immagina un futuro in cui il videogioco sia uno spazio
sicuro, lo immagina come \textbf{spazio condiviso}: ascoltare, parlare e
condividere i videogiochi è ``l'unico modo che abbiamo'' per continuare il
processo di integrazione.

% ---------------------------------------------------------------------
\section{Il punto di vista di chi sviluppa}
\label{sec:c1-1-sviluppatore}

Federico risponde ``come persona che i giochi contribuisce a farli in prima
persona''. Per uno sviluppatore ci sono molti strumenti che obbligano a farsi
domande su questo aspetto. ESRB e PEGI sono ``due ottimi strumenti'' perché negli
anni hanno fatto due cose:
\begin{enumerate}
  \item si sono \textbf{irrigiditi};
  \item hanno adottato \textbf{regole più precise}.
\end{enumerate}
L'idea è che con \textbf{prerequisiti precisi} si ottengono \textbf{risultati
precisi}; senza prerequisiti chiari, il risultato non può essere altrettanto
chiaro.

Per esperienza diretta, assicura che sia dal punto di vista creativo sia da
quello dello sviluppo ci si pone \textbf{quasi da subito} il problema di
\textbf{in quale fascia d'età ricadranno i contenuti del gioco}. Non solo per
ragioni di mercato (la fascia determina un mercato di riferimento, e i giochi
``vengono fatti anche per essere venduti'', oltre che per divertire), ma per
capire \textbf{che tipo di gioco fare} e come affrontare determinati temi. Non
si tratta solo di violenza, di mostrare ``una persona che si fa male'' o di
parolacce: è un discorso ``a tutto tondo'' di \textbf{posizionamento del gioco e
dei suoi contenuti}.

Poiché nella storia del videogioco i temi trattati si sono evoluti e hanno molte
sfaccettature, il gruppo di sviluppo si siede ``al tavolo seriamente'' e si
chiede:
\begin{itemize}
  \item quale sarà il contenuto di questo gioco?
  \item in quale fascia ricadiamo, e perché?
  \item abbiamo bisogno di spingere verso una fascia piuttosto che un'altra?
\end{itemize}
Questi discorsi servono a rendere \textbf{tutto il team consapevole}: ``non ci
possono essere persone che spingono in una direzione e un altro gruppo che
spinge in un'altra''; tutti devono avere le idee chiare. Ciò che uno
sviluppatore può fare, dunque, è \textbf{aprire subito un dialogo} che renda
chiaro a tutti che cosa si sta facendo, in quale direzione si va e come
affrontare i temi in modo che funzionino \emph{sia} per il gioco (altrimenti se
ne snaturerebbe la natura) \emph{sia} per chi il gioco lo fruirà. L'obiettivo è
anche che il prodotto venga \textbf{percepito nella maniera giusta}, cosa ``non
sempre scontata''.

% ---------------------------------------------------------------------
\section{L'evoluzione dei videogiochi italiani}
\label{sec:c1-1-italia}

La moderatrice torna da Davide, che come direttore artistico dell'evento sui
videogiochi italiani osserva l'industria nazionale ``da vicino'': che
evoluzione ha notato nei team italiani? Ci sono caratteristiche che qualche
anno fa non c'erano?

Dopo la battuta su Sanremo (``lo sanno tutti, è il mio grande obiettivo''),
Davide racconta di aver visto in quasi quindici anni, da giornalista e poi
nell'associazione, ``un grande cambiamento'' nell'industria italiana. La sua
storia passa per Milestone e per il genere \textbf{racing}: l'Italia è ``terra
di motori e allenatori di calcio''; sul calcio, nei videogiochi, ``non siamo
andati benissimo quando ci abbiamo provato'', sui motori invece
``fortissimi''. Il settore si è strutturato nel tempo attorno alle aziende
leader\footnote{Il passaggio successivo è in parte una ripetizione meccanica
della trascrizione (``ma insomma, ma insomma\dots''), che non si riporta.}.

Fino a tre--sette anni fa, pensare a un videogioco italiano significava pensare
quasi necessariamente a pochi titoli noti; negli ultimi anni la situazione è
cambiata \textbf{su più livelli}. Segnala un secondo sito, che raccoglie
\textbf{tutti i videogiochi italiani sul mercato}\footnote{L'indirizzo è
trascritto ``www.mesinitali.it'', probabilmente una storpiatura di un dominio
del tipo ``madeinitaly'': non è stato possibile ricostruirlo con certezza.}, oggi
soprattutto quelli pubblicati su Steam (``per il 90\% sostanzialmente''). Da lì
si vede che:
\begin{itemize}
  \item i \textbf{generi coperti} si sono moltiplicati;
  \item la \textbf{quantità} delle produzioni, dalle più indie alle più
        strutturate, è ``aumentata a dismisura'';
  \item è aumentata la \textbf{qualità};
  \item è cresciuto l'approccio all'\textbf{internazionalizzazione}.
\end{itemize}

Su quest'ultimo punto smonta un ``retropensiero'': che un gioco italiano debba
essere in qualche modo \emph{italiano}. Non tutti i giochi francesi hanno
``le baguette e i vestiti a righe'', né bisogna ``mettere il Colosseo in tutti i
videogiochi italiani''. Il videogioco è un'industria \textbf{globale per
definizione}: la maggior parte dei profitti delle aziende non si fa nel proprio
paese, a meno di essere negli Stati Uniti o, in settori specifici, in Giappone o
in Cina. Così, ``forse la cosa migliore che si può dire dei videogiochi italiani
negli ultimi anni è che \textbf{non ti accorgi sempre che sono italiani}'': negli
showcase capita di scoprirlo con meraviglia, ``e non ti dovresti meravigliare''.
L'Italia è partita in ritardo rispetto alle altre nazioni europee, ma oggi
c'è ``finalmente una grande libertà nell'approccio''. ``Non siamo solo una
nazione di racing\footnote{Trascritto ``di resismo'': con ogni probabilità
``di racing''.}: siamo una nazione che fa dei giochi, a volte più riusciti, a
volte meno, come tutte le nazioni.''

% ---------------------------------------------------------------------
\section{Riconoscere un gioco Milestone: l'attenzione al dettaglio}
\label{sec:c1-1-dettaglio}

La domanda a Federico: Milestone è nota per IP consolidate e con un'identità
precisa (\emph{MotoGP} e le altre serie); il nuovo gioco è stato invece ``decisamente
più sperimentale, quasi un azzardo'', riuscito secondo la moderatrice. Se si
togliesse il logo dalla schermata, quale elemento distintivo vorrebbe che facesse
capire al pubblico, ``già dal primo approccio'', che è un gioco Milestone?

Federico premette che non c'è una sola risposta, dipende dal gioco. Ma se
facesse un gioco di corse qualsiasi, vorrebbe che il giocatore, \textbf{alla
prima curva} dopo la partenza, si rendesse conto che è un gioco Milestone per
l'\textbf{attenzione al dettaglio}: in quei ``cinque metri di curva'', il suono;
la possibilità, anche per chi non ha mai preso in mano un gioco del genere, di
affrontare la curva in modo \textbf{soddisfacente}; guardarsi intorno e
riconoscere l'ambiente come ce lo si immagina (``questa curva la conosco, sono
stata a Imola e l'ho vista di persona'').

Il principio generale che ne ricava:
\begin{quote}
  se qualcosa accade in gioco è perché qualcuno dietro ci ha pensato.
\end{quote}
\textbf{Più libertà c'è in un gioco, più persone hanno pensato e permesso quella
libertà}. I giocatori non possono fare cose che programmatori, designer e
artisti non hanno permesso loro di fare; quando succede, ``si capisce subito'',
perché spesso sono \textbf{bug}, cose non previste. Quindi più libertà significa
più studio dietro, e questo avviene grazie all'attenzione al dettaglio e alla
somma di più reparti (design, programmazione, arte).

% ---------------------------------------------------------------------
\section{Il nodo del contatto online}
\label{sec:c1-1-online}

A questo punto la moderatrice ``sottrae per un secondo'' il ruolo di
moderatrice alla presidente\footnote{In realtà è la presidente Zazza a prendere
la parola dicendo di ``sottrarre per un secondo il ruolo di moderatrice'':
l'attribuzione si ricava dal contenuto (cita ``i dati che ho citato prima'').
Il passaggio contiene anche una lunga ripetizione meccanica (``un lavoro che è
un lavoro che è\dots'') che non si riporta.}, che rilancia: gli sviluppatori
cercano di garantire contenuti di qualità e con il messaggio giusto, ma i dati
citati indicano che il problema maggiore oggi è il \textbf{contatto online}.
Lì non c'entra solo chi fa il gioco: i ragazzi possono trovarsi a contatto con
\textbf{terzi malintenzionati}. Il malintenzionato esiste anche fuori dal
digitale, ammette, ``e so che mi rispondereste benissimo'', ma oggi bisogna
occuparsi del mondo online, ``e sarà sempre più così''. Ringrazia la Polizia di
Stato\footnote{``La presenza di Stato'' nella trascrizione: verosimilmente la
Polizia di Stato (Polizia Postale), presente all'evento.} per il lavoro
``impagabile'' di controllo della rete, e chiede agli sviluppatori: come
affrontate il tema quando realizzate un gioco? Con chi ne discutete?

\subsection{Le scelte di infrastruttura}

Federico risponde dalla propria esperienza. Innanzitutto \textbf{non tutti i
giochi} prevedono \textbf{piazze virtuali} in cui molte persone condividono uno
spazio e comunicano. Poi ci sono le \textbf{linee guida dei produttori di
piattaforma} (Sony, Microsoft\footnote{Nella trascrizione i produttori di
console sono chiamati ``platformer'': si intendono i \emph{platform holder}.}).
Le scelte si fanno sull'\textbf{infrastruttura}:
\begin{itemize}
  \item permettere o no ai giocatori di parlarsi;
  \item via \textbf{chat vocale} o \textbf{chat scritta}: le due permettono cose
        diverse, \textbf{sono regolamentate in modo diverso} e richiedono
        implementazioni diverse, più o meno complesse;
  \item in un gioco di corse, \textbf{quanto è necessario} che le persone
        condividano uno spazio in cui possono dirsi quello che vogliono? Con le
        squadre? Che cosa possono dirsi, e per quanto tempo?
\end{itemize}
In una decina d'anni in Milestone non ha mai dovuto sviluppare ``una piazza'':
al massimo ha implementato le linee guida per permettere ai giocatori di
parlare \textbf{solo via chat vocale}, non scritta. Quando si fanno questi
ragionamenti si applicano anche strumenti di \textbf{blocco} e di
\textbf{segnalazione}, perché il giocatore, o chi lo controlla, possa
regolamentare la cosa: se qualcuno mi manda un messaggio che non mi piace,
``lo posso bloccare in qualsiasi momento''\footnote{La trascrizione dice ``lo
posso giocare in un conto momento'': ``lo posso bloccare in qualsiasi
momento''.}.

\subsection{Bloccare non basta}

La presidente obietta: il ragazzo può bloccare, ma \textbf{non è capace di
comprendere il rischio}. Di qui, di nuovo, la necessità di un lavoro in sinergia
tra tutti gli attori; il suo ruolo, ``purtroppo antipatico'', è chiamare tutti a
una responsabilità condivisa, perché ``ancora non siamo arrivati alla tutela,
anzi, forse siamo lontani''. Si dice però contenta degli interventi, che
dimostrano attenzione al tema anche da parte di chi gestisce le piattaforme.

Federico aggiunge un punto: chi fa i giochi ha famiglia e figli e vuole
tutelare ``gli stessi diritti e le stesse possibilità'' che chiede la
presidente; la responsabilità parte anche \textbf{dalle persone che sviluppano}.

\subsection{Rendere sicuro il momento di contatto}

Davide riporta il discorso sul piano generale. Qualunque cosa si faccia, il
\textbf{momento di contatto con il mondo esterno} arriva (``c'è anche nel bar''),
che sia una chat o altro. Come società civile ci si deve interrogare su come
renderlo sicuro, nel senso che:
\begin{itemize}
  \item chi arriva al momento di contatto (il minore, o il genitore che sa che
        il minore vi è esposto) deve avere \textbf{tutti gli strumenti per
        intervenire}, disponibili, chiari, comprensibili;
  \item devono esistere strumenti per aiutare \textbf{istituzioni e
        regolatori} a individuare i momenti di frizione e di attenzione.
\end{itemize}
Secondo lui i produttori di console (Sony, Nintendo, Microsoft) hanno fatto ``un
lavoro incredibile'' nella gestione delle piattaforme, anche online, e nella
semplicità con cui comunicano agli utenti come mettere in sicurezza le cose;
``non è mai abbastanza'', perché la società evolve. Ma l'approccio deve essere
\textbf{più sistemico}: rendere le persone davvero consapevoli parlando,
confrontandosi e \textbf{usando} i prodotti.

Il punto ``vero'', dice riprendendo Federico, è che chi lavora nell'industria è
\textbf{privilegiato}: è cresciuto dentro quei prodotti e sa ``metterci le
mani''. L'esposizione al prodotto e alle sue logiche fa parte di qualsiasi
processo di regolazione. Il videogioco ``si basa sull'interazione'', anche
offline (interagisci con un mondo), e \textbf{l'interazione, se non si esperisce,
non può essere compresa}, né ``preimmaginata''. I videogiochi sono destinati a
restare, fanno parte della cultura di massa, sono sempre più pervasivi ``nel
bene e nel male''; l'unico modo per agire con consapevolezza e responsabilità è
\textbf{conoscerne le logiche e affrontarle tutti insieme}.

\subsection{Non solo console}

Nell'ultimo scambio la presidente concorda sulla consapevolezza: aiuta i
genitori ``e i nonni'' a usare questi strumenti, e cita la propria esperienza
(``io per prima\dots non lo sapevo'') e quella di genitori che ancora oggi non
conoscono strumenti esistenti da anni. Ricorda infine che \textbf{non ci sono
solo le console}: molti giochi passano dai \textbf{telefoni}, con i relativi
messaggi, anche \textbf{pubblicitari}, su cui bisogna lavorare. ``Non basto io,
non bastate voi'': serve un \textbf{tavolo a 360 gradi} con tutti gli attori,
ciascuno responsabile per la sua parte.

% ---------------------------------------------------------------------
\section{Chiusura e premiazione}
\label{sec:c1-1-chiusura}

La moderatrice chiude ammettendo che si è sentita solo ``la punta
dell'iceberg'', con l'auspicio di aver dato spunti di riflessione al pubblico in
sala e online. Approfitta poi della consegna dei premi agli ospiti del festival
(con un microfono ``proprio pirata vero''\footnote{Il senso del passaggio non è
chiaro; potrebbe riferirsi alla forma del premio o a un oggetto scenico.}) per
premiare la presidente Zazza per il suo \textbf{impegno istituzionale}, e per
aver creato ``un dialogo eccezionale a livello istituzionale, cosa rara''.

La presidente aggiunge un'ultima riflessione: \textbf{le istituzioni devono
dialogare}, e il primo problema è spesso proprio la mancanza di apertura al
dialogo. Solo mettendo in campo tutte le competenze e le peculiarità si può
creare consapevolezza e rendere la tutela effettiva: ``l'unione fa la forza,
anche in questo''. La registrazione si chiude tra i saluti e qualche battuta
fuori microfono.

% ---------------------------------------------------------------------
\section*{Punti chiave del capitolo}
\addcontentsline{toc}{section}{Punti chiave del capitolo}

\begin{itemize}
  \item I videogiochi non sono solo intrattenimento: sono un'industria con
        effetti su società, cultura e comunicazione. La \textbf{tutela dei
        minori} deve essere un \textbf{impegno condiviso} di autorità,
        industria, genitori e scuola.
  \item Dati citati: ricerca \textbf{Meter} su 467 bambini (9--15 anni): 70\%
        con almeno una situazione di rischio, solo 10\% pienamente consapevole;
        \textbf{Telefono Azzurro}: nel 2025 il 7,2\% dei casi del 114 riguardava
        abuso, adescamento, ricatto sessuale.
  \item Online e offline non sono mondi separati: la sicurezza degli ambienti
        digitali è come la sicurezza delle piazze. La tutela non è un limite
        all'innovazione, ma \textbf{una misura della sua qualità}.
  \item Richieste all'industria: protezioni per età, controlli parentali
        semplici, comunicazione online più sicura, trasparenza commerciale,
        attenzione anche all'offline. Gli strumenti esistenti vanno resi
        \textbf{conosciuti}.
  \item Il divieto non basta: servono \textbf{alfabetizzazione digitale} e
        pensiero critico. ``Nativo digitale'' non significa ``formato'', e non
        deve diventare un alibi per gli adulti.
  \item L'industria si è \textbf{autoregolata} (ESRB, poi PEGI), per
        autocoscienza e per ragioni commerciali; l'autoregolazione è l'inizio di
        un dialogo. Strumenti utili: \textbf{tuttosuivideogiochi.it}, i patti
        di famiglia sulle regole condivise.
  \item Per chi sviluppa, la \textbf{fascia d'età} si discute da subito e
        coinvolge tutto il team: non è solo mercato, ma posizionamento dei temi
        e dei contenuti.
  \item L'industria italiana è cresciuta in quantità, qualità, generi e
        internazionalizzazione: il segno migliore è che spesso \textbf{non ti
        accorgi che un gioco è italiano}.
  \item ``Se qualcosa accade in gioco è perché qualcuno ci ha pensato'': più
        libertà significa più progettazione dietro.
  \item Sul contatto online si decide a livello di \textbf{infrastruttura}
        (chat vocale o scritta, se e quanto serve), con blocchi e segnalazioni;
        ma il minore non sempre comprende il rischio, e l'interazione si capisce
        solo \textbf{facendone esperienza}.
\end{itemize}
